\chapter{Documentation API détaillée}

Cette annexe présente une synthèse technique des principales API exposées par le backend Laravel de \textbf{SomaScan}. Toutes les routes applicatives sont préfixées par \codepath{/api} et utilisent le format JSON. Les applications mobile Flutter et web React consomment ces endpoints afin d'assurer l'authentification, la gestion des camions, le suivi des trajets, la consultation des logs et l'accès aux indicateurs opérationnels.

\section{Principes généraux}

Les requêtes vers l'API utilisent généralement les en-têtes suivants :

\begin{lstlisting}[language={},caption={En-têtes HTTP utilisés par les clients API}]
Accept: application/json
Content-Type: application/json
Authorization: Bearer <token>
\end{lstlisting}

Les réponses applicatives suivent une structure normalisée. Cette convention permet aux applications clientes de traiter les succès, les erreurs fonctionnelles et les erreurs de validation de manière cohérente.

\begin{lstlisting}[language={},caption={Structure générale d'une réponse API}]
{
  "success": true,
  "data": {},
  "message": "Optional message",
  "errors": {}
}
\end{lstlisting}

L'authentification repose sur des jetons Laravel Sanctum. Les endpoints publics sont accessibles sans jeton, tandis que les endpoints protégés exigent un jeton valide dans l'en-tête \codepath{Authorization}.

\section{Authentification}

\small
\begin{longtable}{|p{2.3cm}|p{4.2cm}|p{2.8cm}|p{5.1cm}|}
\hline
\textbf{Méthode} & \textbf{Endpoint} & \textbf{Accès} & \textbf{Rôle} \\
\hline
\endhead

POST & \codepath{/api/login} & Public & Authentifier un utilisateur et générer un jeton Sanctum. \\
\hline
POST & \codepath{/api/logout} & Authentifié & Révoquer le jeton courant de l'utilisateur connecté. \\
\hline
GET & \codepath{/api/me} & Authentifié & Retourner le profil de l'utilisateur connecté. \\
\hline
\caption{Endpoints d'authentification}
\label{tab:api-authentication}
\end{longtable}
\normalsize

\begin{lstlisting}[language={},caption={Exemple de requête de connexion}]
{
  "email": "user@example.com",
  "password": "secretpassword",
  "device_name": "android-tablet-1"
}
\end{lstlisting}

\begin{lstlisting}[language={},caption={Exemple de réponse de connexion}]
{
  "token": "1|abc...",
  "user": {
    "id": 1,
    "name": "Admin",
    "role": "ADMIN",
    "location": null
  },
  "expires_at": null
}
\end{lstlisting}

\section{Endpoints des opérateurs terrain}

Ces endpoints sont destinés aux opérateurs affectés à l'usine ou au port. Le endpoint central est \codepath{/api/scan}, car il déclenche l'évolution du cycle de vie d'un trajet selon l'état courant du camion et la position de l'opérateur.

\small
\begin{longtable}{|p{2.3cm}|p{4.4cm}|p{3.3cm}|p{4.4cm}|}
\hline
\textbf{Méthode} & \textbf{Endpoint} & \textbf{Rôles} & \textbf{Description} \\
\hline
\endhead

POST & \codepath{/api/scan} & \codepath{COMPANY_OPERATOR}, \codepath{PORT_OPERATOR} & Soumettre un scan QR et mettre à jour automatiquement le statut du trajet. \\
\hline
GET & \codepath{/api/operator/last-scans} & \codepath{COMPANY_OPERATOR}, \codepath{PORT_OPERATOR} & Consulter les derniers scans effectués par l'opérateur connecté. \\
\hline
GET & \codepath{/api/trucks/{id}/basic} & Tous les rôles authentifiés & Récupérer les informations légères d'un camion après un scan. \\
\hline
\caption{Endpoints utilisés par les opérateurs terrain}
\label{tab:api-operator-endpoints}
\end{longtable}
\normalsize

\begin{lstlisting}[language={},caption={Exemple de requête de scan QR}]
{
  "qr_code": "SOMASTEEL-TRUCK-24452-A-54",
  "device_time": "2026-06-22T10:00:00Z",
  "device_id": "device-uuid-123"
}
\end{lstlisting}

\begin{lstlisting}[language={},caption={Exemple de réponse après scan réussi}]
{
  "status": "SUCCESS",
  "message": "Scan successful",
  "current_step": "STARTED",
  "next_expected_step": "ARRIVED_PORT",
  "is_locked": true,
  "trip_summary": {
    "trip_id": 42,
    "truck_id": 5,
    "status": "STARTED",
    "truck": {
      "id": 5,
      "registration_number": "24452 A 54",
      "driver_name": "Nom chauffeur"
    },
    "action": "START",
    "timestamps": {
      "started_at": "2026-06-22T10:00:00Z",
      "arrived_port_at": null,
      "left_port_at": null,
      "completed_at": null
    }
  }
}
\end{lstlisting}

Les principales erreurs retournées par le scan sont les suivantes : camion introuvable, camion inactif, rôle ou localisation incompatible avec l'action demandée, trajet déjà terminé ou flux de scan non configuré.

\section{Gestion des camions}

\small
\begin{longtable}{|p{2.3cm}|p{4.8cm}|p{2.5cm}|p{4.8cm}|}
\hline
\textbf{Méthode} & \textbf{Endpoint} & \textbf{Rôle} & \textbf{Description} \\
\hline
\endhead

GET & \codepath{/api/trucks} & ADMIN & Lister les camions avec pagination et filtres. \\
\hline
POST & \codepath{/api/trucks} & ADMIN & Créer un camion et générer son QR code si nécessaire. \\
\hline
GET & \codepath{/api/trucks/{id}} & ADMIN & Consulter les détails d'un camion. \\
\hline
PUT & \codepath{/api/trucks/{id}} & ADMIN & Modifier les informations d'un camion. \\
\hline
DELETE & \codepath{/api/trucks/{id}} & ADMIN & Supprimer un camion et les données associées. \\
\hline
POST & \codepath{/api/trucks/{id}/generate-qr} & ADMIN & Régénérer le code QR à partir de l'immatriculation. \\
\hline
PATCH & \codepath{/api/trucks/{id}/activate} & ADMIN & Activer un camion. \\
\hline
PATCH & \codepath{/api/trucks/{id}/deactivate} & ADMIN & Désactiver un camion. \\
\hline
POST & \codepath{/api/trucks/bulk-activate} & ADMIN & Activer plusieurs camions en une seule opération. \\
\hline
POST & \codepath{/api/trucks/bulk-deactivate} & ADMIN & Désactiver plusieurs camions en une seule opération. \\
\hline
GET & \codepath{/api/trucks/{id}/trips} & ADMIN & Consulter les trajets associés à un camion. \\
\hline
GET & \codepath{/api/trucks/search} & ADMIN & Rechercher des camions par immatriculation ou chauffeur. \\
\hline
GET & \codepath{/api/trucks/stats} & ADMIN & Obtenir les statistiques globales du parc. \\
\hline
\caption{Endpoints de gestion des camions}
\label{tab:api-trucks}
\end{longtable}
\normalsize

\section{Gestion des trajets}

\small
\begin{longtable}{|p{2.3cm}|p{4.8cm}|p{2.5cm}|p{4.8cm}|}
\hline
\textbf{Méthode} & \textbf{Endpoint} & \textbf{Rôle} & \textbf{Description} \\
\hline
\endhead

GET & \codepath{/api/trips} & ADMIN & Lister les trajets avec pagination et filtres. \\
\hline
GET & \codepath{/api/trips/active} & ADMIN & Lister les trajets actifs. \\
\hline
GET & \codepath{/api/trips/history} & ADMIN & Lister l'historique des trajets terminés. \\
\hline
GET & \codepath{/api/trips/{id}} & ADMIN & Consulter les détails d'un trajet. \\
\hline
GET & \codepath{/api/trips/{id}/logs} & ADMIN & Consulter les logs de scan d'un trajet. \\
\hline
GET & \codepath{/api/trips/{id}/timeline} & ADMIN & Obtenir la chronologie des événements d'un trajet. \\
\hline
PATCH & \codepath{/api/trips/{id}/cancel} & ADMIN & Annuler un trajet actif avec une note optionnelle. \\
\hline
PATCH & \codepath{/api/trips/{id}/notes} & ADMIN & Modifier les notes associées à un trajet. \\
\hline
DELETE & \codepath{/api/trips/{id}} & ADMIN & Supprimer définitivement un trajet. \\
\hline
GET & \codepath{/api/trips/search} & ADMIN & Rechercher des trajets avec filtres. \\
\hline
GET & \codepath{/api/trips/stats} & ADMIN & Obtenir les statistiques des trajets. \\
\hline
GET & \codepath{/api/trips/calendar} & ADMIN & Obtenir une synthèse journalière pour la vue calendrier. \\
\hline
GET & \codepath{/api/trips/by-day} & ADMIN & Obtenir les trajets d'une journée donnée. \\
\hline
\caption{Endpoints de gestion des trajets}
\label{tab:api-trips}
\end{longtable}
\normalsize

\begin{lstlisting}[language={},caption={Exemple de réponse de la vue calendrier}]
{
  "data": [
    {
      "day": "2026-06-22",
      "start_at": "2026-06-22T07:00:00+01:00",
      "end_at": "2026-06-23T06:59:59+01:00",
      "total": 15,
      "active": 3,
      "completed": 12,
      "by_status": {
        "STARTED": 1,
        "ARRIVED_PORT": 1,
        "LEFT_PORT": 1,
        "COMPLETED": 12
      },
      "trips": []
    }
  ],
  "meta": {
    "from": "2026-06-22",
    "to": "2026-06-30",
    "day_start": "07:00",
    "timezone": "Africa/Casablanca"
  }
}
\end{lstlisting}

\section{Gestion des utilisateurs}

\small
\begin{longtable}{|p{2.3cm}|p{4.8cm}|p{2.5cm}|p{4.8cm}|}
\hline
\textbf{Méthode} & \textbf{Endpoint} & \textbf{Rôle} & \textbf{Description} \\
\hline
\endhead

GET & \codepath{/api/users} & ADMIN & Lister les utilisateurs avec filtres par rôle ou localisation. \\
\hline
POST & \codepath{/api/users} & ADMIN & Créer un utilisateur. \\
\hline
GET & \codepath{/api/users/{id}} & ADMIN & Consulter les détails d'un utilisateur. \\
\hline
PUT & \codepath{/api/users/{id}} & ADMIN & Modifier les informations d'un utilisateur. \\
\hline
PATCH & \codepath{/api/users/{id}/reset-password} & ADMIN & Réinitialiser le mot de passe d'un utilisateur. \\
\hline
GET & \codepath{/api/users/{id}/activity} & ADMIN & Consulter l'activité d'un utilisateur. \\
\hline
GET & \codepath{/api/users/search} & ADMIN & Rechercher un utilisateur par nom, email ou rôle. \\
\hline
GET & \codepath{/api/users/stats} & ADMIN & Obtenir les statistiques des utilisateurs. \\
\hline
\caption{Endpoints de gestion des utilisateurs}
\label{tab:api-users}
\end{longtable}
\normalsize

La route \codepath{DELETE /api/users/{id}} est enregistrée dans la couche de routage, mais la méthode de suppression correspondante n'est pas encore implémentée dans le contrôleur. Elle ne doit donc pas être considérée comme une fonctionnalité opérationnelle dans l'état actuel.

\section{Logs de scan, rapports et configuration}

\small
\begin{longtable}{|p{2.3cm}|p{5cm}|p{2.5cm}|p{4.6cm}|}
\hline
\textbf{Méthode} & \textbf{Endpoint} & \textbf{Rôle} & \textbf{Description} \\
\hline
\endhead

GET & \codepath{/api/scan-logs} & ADMIN & Lister et filtrer l'ensemble des scans effectués. \\
\hline
GET & \codepath{/api/reports/summary} & ADMIN & Obtenir une synthèse analytique des trajets. \\
\hline
GET & \codepath{/api/reports/truck/{id}} & ADMIN & Obtenir un rapport analytique pour un camion donné. \\
\hline
GET & \codepath{/api/reports/export} & ADMIN & Télécharger un export Excel des trajets. \\
\hline
GET & \codepath{/api/scan-flow} & ADMIN & Consulter le flux de statuts actif. \\
\hline
PUT & \codepath{/api/scan-flow} & ADMIN & Modifier la séquence logique du flux de scan. \\
\hline
\caption{Endpoints d'administration complémentaires}
\label{tab:api-admin-extra}
\end{longtable}
\normalsize

Les endpoints \codepath{/api/reports/durations} et \codepath{/api/reports/delays} sont présents dans la couche de routage mais leurs méthodes dédiées ne sont pas encore implémentées. Ils ne sont donc pas retenus comme fonctionnalités finalisées.

\section{Endpoint de configuration initiale}

Le endpoint \codepath{GET /api/setup} est public et permet d'exécuter les migrations et seeders lors de l'initialisation. Il doit être considéré comme un endpoint technique de bootstrap, à limiter fortement ou désactiver en production afin d'éviter toute action non contrôlée sur la base de données.
